%% Response to Reviewers — generated by make_manuscript.py from outputs/revision.
%% Set \internalfalse to hide author-only notes before sending.
\documentclass[11pt]{article}
\usepackage[margin=1in]{geometry}\usepackage[T1]{fontenc}\usepackage[utf8]{inputenc}
\usepackage{amsmath,amssymb,booktabs,array,xcolor,enumitem,listings,graphicx}
\usepackage{tcolorbox}\tcbuselibrary{breakable,skins}
\usepackage[colorlinks=true,linkcolor=blue,urlcolor=blue]{hyperref}
\graphicspath{{./}{./figures/}}
\newif\ifinternal \internaltrue
\definecolor{cmtbg}{HTML}{F2F4F7}\definecolor{cmtfr}{HTML}{B0B7C3}\definecolor{bblue}{HTML}{1565C0}\definecolor{bred}{HTML}{C62828}\definecolor{notebg}{HTML}{FFF8E1}\definecolor{notefr}{HTML}{C8A200}
\newtcolorbox{rcomment}[1]{breakable,colback=cmtbg,colframe=cmtfr,boxrule=0.5pt,arc=2pt,left=7pt,right=7pt,top=6pt,bottom=6pt,title={#1},fonttitle=\bfseries\small,coltitle=black,colbacktitle=cmtfr!45}
\newcommand{\response}{\par\medskip\noindent\textbf{\color{bblue}Response.}\ }
\newcommand{\actions}{\par\medskip\noindent\textbf{Changes to the manuscript.}}
\newcommand{\pending}[1]{\textcolor{bred}{\textbf{[TO FILL: #1]}}}
\newcommand{\internalnote}[1]{\ifinternal\begin{tcolorbox}[breakable,colback=notebg,colframe=notefr,boxrule=0.5pt,arc=2pt]\footnotesize\textbf{Internal note --- remove before submission.}\ #1\end{tcolorbox}\fi}
\lstset{basicstyle=\ttfamily\footnotesize,breaklines=true,frame=single,rulecolor=\color{cmtfr},backgroundcolor=\color{cmtbg}}
\setlist[itemize]{leftmargin=1.4em,itemsep=2pt,topsep=3pt}
\title{\vspace{-2em}\textbf{Response to Reviewers}\\[0.4em]\large Manuscript a69523d8-f3cf-41cd-b7ab-1b70f2f037b1 (\emph{Scientific Reports})\\[0.2em]\normalsize revised as: \emph{Data-Efficient Nucleus Segmentation of Mammary Epithelial Cells in Confocal Stacks: What a Leakage-Free, Budget-Matched Evaluation Reveals About Handcrafted Additions}}
\author{Varun Kumar Dasoju, Qingsu Cheng, Zeyun Yu}\date{\today}
\begin{document}\maketitle\thispagestyle{empty}
\begin{center}\textcolor{bred}{\textbf{INTERNAL ARCHIVE --- DO NOT SUBMIT TO A NEW JOURNAL}}\end{center}
\noindent This response addresses reviews of an earlier submission. Several responses below describe the state of that revision at the time and are retained only as an audit trail. The next-journal manuscript and current methodological limitations take precedence.

\section*{Preface}
We thank the Editorial Board Member and both reviewers. The criticisms were well founded, and in acting on them we audited our dataset, code and released checkpoint directly. That audit confirmed the reviewers' central concern about train/test independence and uncovered two further defects in our own pipeline that neither reviewer could have seen. We have therefore rebuilt the evaluation from the data split upward and recomputed every number. Rather than defend the original figures we report the corrected ones, which are lower, together with the reasons.

\section{Summary of principal changes}
\begin{enumerate}[nosep]
\item \textbf{Volume-wise split.} The original split was slice-level: @@LEAK_ORIG_95@@\% of test images had a near-duplicate (NCC $>$0.95) in the training set (median max-NCC @@LEAK_ORIG_MED@@). Slices are now partitioned by acquisition field (@@V_TRAIN@@/@@V_VAL@@/@@V_TEST@@ fields; @@N_TRAIN@@/@@N_VAL@@/@@N_TEST@@ slices); maximum cross-split NCC is @@LEAK_NEW_MAX@@. Every model was retrained.
\item \textbf{Dataset statistics corrected.} The ``60\% empty images'' in the original Table~1 was an encoding artefact (0/1-stored masks read at a 127 threshold). All @@N_TOTAL@@ valid slices contain nuclei; two corrupt label files were found and excluded.
\item \textbf{Objective corrected.} Algorithm~1 now describes exactly the loss that trained the released model (stabilised Dice+BCE). The Tversky and boundary terms it previously claimed are reported as ablations, with the correct recall-biased parameterisation.
\item \textbf{Identical budgets.} One harness, one 30-epoch protocol, one preprocessing path for every model, including the two baselines previously omitted; three seeds for the proposed model.
\item \textbf{Expanded ablations and baselines.} 11 single-change ablations, 4 alternative edge detectors, 6 Gabor sensitivity runs, a 6-encoder ladder, a vanilla anchor, and Cellpose-SAM zero-shot and fine-tuned.
\item \textbf{Statistics.} BCa bootstrap CIs, paired Wilcoxon tests with Holm correction, Cliff's $\delta$; boundary metrics (HD95, ASSD, NSD$_2$) throughout.
\item \textbf{Claims narrowed.} Clinical-grade, ``boundary-aware'' and generalisation claims removed; two implementation bugs disclosed (Gabor normalisation, sampler threshold).
\item \textbf{Reproducibility.} Code, split file, seeds, per-image metrics and weights released; figures regenerated from code; AI-use statement added.
\item \textbf{The paper's claim has changed.} Under the corrected protocol none of the four additions we previously called essential produces an effect larger than seed-to-seed variation (Gabor $\Delta$=@@DELTA:abl_no_gabor@@, SCSE @@DELTA:abl_no_scse@@, sampling @@DELTA:abl_uniform_sampling@@, stabilised loss @@DELTA:abl_loss_dicebce@@; seed SD @@PROP_SEED_SD@@), and a vanilla pretrained EfficientNet-B7 UNet++ with plain Dice+BCE (@@D:vanilla_effb7_unetpp@@) is the strongest configuration. The manuscript is reframed around the audit, the corrected recipe, and the specific defects that manufactured the original gaps. We considered this more useful than a weaker restatement of the original claims.
\end{enumerate}

\section{Self-audit findings}\label{sec:audit}
\subsection*{Finding 1 --- the split leaked}
Each 2D image is one $z$-plane of a confocal stack of one organoid. Assigning images to splits individually put adjacent planes of the same organoid on both sides. Normalised cross-correlation of $128^2$ thumbnails, test vs.\ all training images: median max-NCC @@LEAK_ORIG_MED@@; @@LEAK_ORIG_95@@\% above 0.95; @@LEAK_ORIG_90@@\% above 0.90; maximum @@LEAK_ORIG_MAX@@. Visual inspection confirmed identical organoids at adjacent depths (Fig.~1 of the revised manuscript). Under the new field-wise split: median @@LEAK_NEW_MED@@, maximum @@LEAK_NEW_MAX@@, none above 0.80. Re-evaluating the originally submitted checkpoint with its original preprocessing on the original test set gives Dice @@REL_ALL512@@ at $512^2$ (the paper reported 0.951 on a 100-image subset); on the @@REL_N_LEAK@@ test images with a training near-duplicate it scores @@REL_LEAK@@ (native resolution) and on the @@REL_N_NOLEAK@@ without, @@REL_NOLEAK@@.

\subsection*{Finding 2 --- ``60\% empty'' was an encoding artefact}
Masks are stored in two encodings: 347 as 0/255 and 508 as 0/1. Our dataset statistics thresholded at 127, reading every 0/1 mask as empty. Visual inspection and mask--image correlation ($r$=0.82) show they are genuine annotations (mean foreground 8.3\%). Two further files contained 16-bit image data instead of a mask and were excluded. The training loader binarised both encodings correctly, so the model was trained on valid labels---but the \emph{sampler} used the same 127 threshold and therefore treated 508 nucleus-containing images as background (weight 0.3, $8.9\times$ under-sampled). Both are fixed and disclosed.

\subsection*{Finding 3 --- Algorithm 1 did not match the trained model}
The released checkpoint was trained with $0.7\,\mathcal{L}_{\rm Dice}+0.3\,\mathcal{L}_{\rm BCE}^{\rm adaptive}$: no Tversky term, no boundary term. The manuscript's Algorithm~1 described a different objective. We now document the objective actually used and evaluate the Tversky+boundary variant as an ablation.

\subsection*{Headline result under the corrected protocol}
Proposed configuration on @@N_TEST@@ slices from @@V_TEST@@ unseen fields: Dice @@PROP_DICE@@ [@@PROP_DICE_LO@@, @@PROP_DICE_HI@@], IoU @@PROP_IOU@@, precision @@PROP_P@@, recall @@PROP_R@@, HD95 @@PROP_HD95@@\,$\mu$m, ASSD @@PROP_ASSD@@\,$\mu$m; @@N_SEEDS@@ seeds @@PROP_SEED_MEAN@@$\pm$@@PROP_SEED_SD@@. Best budget-matched baseline (@@BEST_BASE@@) @@BEST_BASE_DICE@@; zero-shot Cellpose-SAM @@CELLPOSE_ZEROSHOT_DICE@@.

\section{Editor}
\begin{rcomment}{Editorial Board Member}Based on Reviewer 1 there are fundamental issues that cannot be addressed. Reviewer 2: ablations insufficient; add statistical tests; more dataset detail; more SOTA baselines.\end{rcomment}
\response We agree the issues were fundamental as submitted. They were addressed by rebuilding the evaluation (Section~\ref{sec:audit}) rather than by editing text. Ablations: R1.6, R1.8, R1.17, R2.6. Statistics: R1.12. Dataset: R1.1, R1.2, R1.18, R2.8. Baselines: R1.7, R1.20, R2.9.

\section{Reviewer 1}
\subsection*{1. Empty-image inclusion and Dice on empty masks}
\begin{rcomment}{R1.1}Empty-image inclusion may inflate Dice. Report metrics separately for nucleus-containing and empty images and define Dice handling for empty masks.\end{rcomment}
\response The reviewer's concern was reasonable given what we wrote, but on audit the premise turned out to be wrong in our own favour and our own fault: the dataset contains no empty images (Finding~2). The ``60\%'' was an encoding artefact in our statistics script. Table~1 has been recomputed from the files; the empty-mask convention is nevertheless stated in Methods (Dice undefined for $|P|=|T|=0$; specificity and false-positive pixels would be reported for such images), and a pooled micro-averaged Dice (@@PROP_MICRO@@ for the proposed model) is reported alongside the per-image mean.
\actions Methods §Dataset (encodings, exclusions, foreground range @@FG_MIN@@--@@FG_MAX@@\%); Methods §Evaluation (convention, micro-average); Table~1 replaced.

\subsection*{2. Train/test independence}
\begin{rcomment}{R1.2}Specify whether splits are volume-wise and ensure no source-volume or adjacent-slice leakage.\end{rcomment}
\response Confirmed and corrected (Finding~1). Splits are now by acquisition field; the assignment file is released; the leakage audit is reported for both splits (revised Fig.~1); every model was retrained on the new split.
\actions Methods §Dataset (provenance recovery, split, audit); Results §``Dataset provenance''; Table~2 (released-checkpoint audit); all result tables recomputed.

\subsection*{3. Reproducibility}
\begin{rcomment}{R1.3}Publicly release code and a representative data sample.\end{rcomment}
\response Agreed. The complete harness (\texttt{train\_revision.py}, \texttt{stats.py}, \texttt{figures.py}, \texttt{make\_manuscript.py}), the split file, seeds, per-image metric CSVs and the proposed model's weights are released; every number in the paper is regenerated from them by one command. A representative annotated sample is released under a data-use agreement.
\actions Data availability rewritten. \internalnote{Create the public repo and Zenodo DOI and fill the three \texttt{\textbackslash pending} markers in Data availability; decide the sample size and licence with Q.C.}

\subsection*{4. Tversky $\alpha$/$\beta$}
\begin{rcomment}{R1.4}Tversky $\alpha$ and $\beta$ interpretation is inconsistent with recall bias.\end{rcomment}
\response Correct. With $T=TP/(TP+\alpha FP+\beta FN)$, recall bias requires $\beta>\alpha$; the manuscript stated $\alpha$=0.7, $\beta$=0.3 and then gave a self-contradictory explanation. The deeper problem (Finding~3) is that the released model contained no Tversky term. The revision documents the objective actually used (stabilised Dice+BCE) and reports a Tversky+boundary objective with $\alpha$=0.3, $\beta$=0.7---genuinely recall-biased---as an ablation: Dice @@D:abl_loss_tversky_boundary@@ vs.\ @@PROP_DICE@@ ($p$=@@P:abl_loss_tversky_boundary@@).
\actions Algorithm~1 rewritten; Methods §Objective; Table~4.

\subsection*{5. Boundary-aware claims}
\begin{rcomment}{R1.5}Rename the loss or add an explicit boundary term and boundary-specific evaluation.\end{rcomment}
\response Both. ``Boundary-aware'' is removed from the title, abstract and claims because the released model had no boundary term. Boundary metrics (HD95, ASSD, NSD$_2$) are now reported for every model at native resolution; the proposed model's ASSD is @@PROP_ASSD@@\,$\mu$m and NSD$_2$ @@PROP_NSD@@. An explicit morphological-gradient boundary term is evaluated as an ablation (previous item).
\actions Title/abstract; Methods §Evaluation; all tables.

\subsection*{6. Gabor sensitivity and alternative detectors}
\begin{rcomment}{R1.6}Provide ablations over orientations, scales, sigma, and alternative edge detectors.\end{rcomment}
\response Done (Table~5, Table~4). Orientations 4/16: @@D:gabor_orient4@@/@@D:gabor_orient16@@; scales 2/5: @@D:gabor_scale2@@/@@D:gabor_scale5@@; $\sigma$=2/10: @@D:gabor_sigma2@@/@@D:gabor_sigma10@@ (default: @@PROP_DICE@@). Alternative channels: Sobel @@D:edge_sobel@@, Scharr @@D:edge_scharr@@, LoG @@D:edge_log@@, Canny @@D:edge_canny@@; no edge channel @@D:abl_no_gabor@@ ($p$=@@P:abl_no_gabor@@). Two corrections: the wavelengths are $\{20,6.67,4\}$\,px, not $\{4,10,20\}$; and the kernels were normalised by their sum, which is near zero for 8 of the 24 kernels---we now use L1 normalisation and say so. Fig.~5 now shows the real kernels and responses.
\actions Methods §Gabor; Tables~4--5; Fig.~5.

\subsection*{7. Baseline budgets}
\begin{rcomment}{R1.7}Train all baselines with identical epochs, schedules, and stopping criteria.\end{rcomment}
\response The reviewer was right: baselines had 150 epochs and the proposed model 200. All models are now trained by one script under one protocol (30 epochs, OneCycleLR, AdamW $3\times10^{-4}$, batch 8, identical augmentation and selection rule) and the best epoch and wall-clock are reported. Under this protocol the baselines score @@WORST_BASE_DICE@@--@@BEST_BASE_DICE@@ rather than 0.47--0.76. We also tested a suspected cause---the original loader truncated 16-bit images to their high byte (maximum $\sim$11/255)---and found it costs only $\sim$0.01 Dice (U-Net R34 @@D:legacy_prep_unet_r34@@ vs.\ @@D:base_unet_r34@@; proposed @@D:legacy_prep_proposed@@ vs.\ @@PROP_DICE@@). It does not explain the original 0.47--0.76, whose cause we cannot reconstruct with confidence; the unequal budget and the leaking split are the demonstrable factors, and we say so in the paper.
\actions Methods §Training protocol; Table~3.

\subsection*{8. Weighted sampling}
\begin{rcomment}{R1.8}Provide sampling formulas, probabilities, and ablations.\end{rcomment}
\response The formula is now in Methods. Auditing it revealed that the original sampler's threshold mis-read the 0/1-encoded masks (Finding~2), so it did not implement the described prioritisation. With the corrected sampler, uniform sampling gives @@D:abl_uniform_sampling@@ vs.\ @@PROP_DICE@@ ($p$=@@P:abl_uniform_sampling@@).
\actions Methods §Sampling (formula, probabilities, bug disclosure); Table~4.

\subsection*{9. Zero-positive batches}
\begin{rcomment}{R1.9}Specify safe handling for zero-positive batches and verify numerical stability.\end{rcomment}
\response The guard ($r_{\rm pos}=0\Rightarrow w_{\rm pos}=1$), the logit clamp and a non-finite fallback existed in code but not in Algorithm~1; all three are now shown. Because every image contains nuclei, zero-positive batches do not occur in practice; the harness counts non-finite loss events per run (proposed model: @@PROP_NONFINITE@@).
\actions Algorithm~1; Results §Training behaviour.

\subsection*{10. Annotation quality}
\begin{rcomment}{R1.10}Report annotator count, protocol, agreement metrics, and boundary uncertainty.\end{rcomment}
\response The protocol is now described (single trained annotator, expert review, tool, binary masks). The unmeasured ``$<$3 px'' claim has been removed. We have not yet run a formal multi-annotator study and say so as the most important missing validation. However, 4-fold cross-validation over every field (new Table~7) measures the effect indirectly and decisively: the two annotation sessions (identifiable by mask encoding) used different boundary conventions---@@CV_TIGHT_N@@ fields give recall $\geq$0.98 with precision 0.72--0.89 (masks tighter than the model), @@CV_WIDE_N@@ fields give precision $\geq$0.96 with recall 0.73--0.83 (masks wider)---and cross-convention testing costs 0.05--0.10 Dice. The manuscript now reports this, states that annotation convention sets the ceiling, and recommends fixing one convention before annotating.
\actions Methods §Dataset; Discussion §Limitations. \internalnote{An inter-annotator study needs a second annotator on $\geq$30 slices; if it can be done before resubmission, add pairwise Dice, boundary distance and $\kappa$.}

\subsection*{11. Clinical-grade claims}
\begin{rcomment}{R1.11}Remove clinical-grade deployment claims or provide clinical validation.\end{rcomment}
\response Removed entirely: abstract, introduction, Discussion (including the ``clinical threshold'' line and the surgical-margin analogy, which does not apply to cultured cells), and ``clinical validation'' in the contributions statement. The work is framed as research-grade morphometry.

\subsection*{12. Confidence intervals and robust tests}
\begin{rcomment}{R1.12}Report bootstrap confidence intervals and appropriate paired or nonparametric tests.\end{rcomment}
\response Every metric carries a 95\% BCa bootstrap CI (10\,000 resamples over test images). Comparisons use the paired Wilcoxon signed-rank test on per-image Dice with Holm--Bonferroni correction across all comparisons, plus Cliff's $\delta$. The $t$-test has been dropped as inappropriate for bounded, skewed per-image scores.
\actions Methods §Evaluation; every table.

\subsection*{13. Figure labels and AI visuals}
\begin{rcomment}{R1.13}Correct ``Quantum'' labels, provide actual Gabor examples, disclose AI generation.\end{rcomment}
\response ``Quantum'' (a legacy class name) is removed from all figures, code and text. The former illustrative panels were generated by a language model and are replaced by Fig.~5, produced by the released code from a held-out test slice (kernels, individual responses, the max-response channel, Sobel for comparison, ground truth). An AI-use statement is added to Methods.

\subsection*{14. External validation}
\begin{rcomment}{R1.14}Evaluate on an independent external dataset or restrict generalization claims.\end{rcomment}
\response We restrict the claims. All data are MCF10A organoids from one laboratory and instrument; the Discussion now states this as a limitation and makes no generalisation claim. \internalnote{Internal domain shift (184A1, 184D lines; 1,531 volumes) is available but has no 2D ground truth; a small annotation effort there would answer this comment directly.}

\subsection*{15. Test-time augmentation}
\begin{rcomment}{R1.15}State whether reported metrics use augmentation and apply it identically to baselines.\end{rcomment}
\response The ``7-fold TTA'' statement was wrong (the code implemented 3 transforms and the evaluation script did not call it). Every model is now evaluated both without TTA (all headline numbers) and with identical 3-fold flip TTA; for the proposed model TTA gives @@PROP_TTA_DICE@@ vs.\ @@PROP_DICE@@.
\actions Methods §Evaluation; per-run \texttt{result.json}.

\subsection*{16. Overfitting safeguards}
\begin{rcomment}{R1.16}Provide learning curves, cross-validation, and regularization details.\end{rcomment}
\response Validation is on held-out volumes; training/validation curves for @@N_SEEDS@@ seeds are shown (Fig.~7); @@CV_N@@-fold leave-fields-out cross-validation over all @@N_VOL@@ fields gives pooled Dice @@CV_POOLED@@ (fold means @@CV_MEAN@@$\pm$@@CV_SD@@; @@CV_ROWS@@), lower than the hold-out @@PROP_DICE@@ because field difficulty varies---both are reported; all regularisers are listed with values in Methods §Training protocol; the encoder ladder (Table~6) addresses capacity directly.

\subsection*{17. Isolating EfficientNet-B7}
\begin{rcomment}{R1.17}Include an EfficientNet-B7 baseline without proposed additions and report encoder ablations.\end{rcomment}
\response Both added, and this experiment turned out to be decisive. Vanilla EfficientNet-B7 UNet++ (RGB, no SCSE, uniform sampling, plain Dice+BCE) scores @@D:vanilla_effb7_unetpp@@ vs.\ @@PROP_DICE@@ for the full configuration ($p$=@@P:vanilla_effb7_unetpp@@): the additions do not help. The paper now says so. Encoder ladder with everything else fixed: ResNet-34 @@D:ladder_resnet34@@, ResNet-50 @@D:ladder_resnet50@@, EfficientNet-B0 @@D:ladder_effb0@@, B3 @@D:ladder_effb3@@, B5 @@D:ladder_effb5@@, B7 @@PROP_DICE@@.
\actions Table~4 (last row); Table~6; Fig.~6.

\subsection*{18. ROI extraction}
\begin{rcomment}{R1.18}Describe crop generation, number of crops per source image, and split assignment.\end{rcomment}
\response Each image is a full acquisition frame (one per $z$-plane, no sub-crops). The ROI stage of the original pipeline---which was applied in training but not at test time, and with 30-px padding rather than the stated 50---has been removed entirely so that training and evaluation share one preprocessing path. Split assignment is by acquisition field and is released.
\actions Methods §Dataset, §Preprocessing.

\subsection*{19. IoU inconsistency}
\begin{rcomment}{R1.19}Reconcile the abstract IoU value with the results.\end{rcomment}
\response Corrected. Every number in the revised manuscript (abstract, tables, text) is inserted programmatically from a single results file, so they cannot diverge. The revised IoU is @@PROP_IOU@@.

\subsection*{20. Omitted baselines}
\begin{rcomment}{R1.20}Include final metrics or failure analyses for all attempted baselines.\end{rcomment}
\response MA-Net and UNet++ (ResNet-50) are included under the shared protocol: @@D:base_manet_r50@@ and @@D:base_unetpp_r50@@. Under equal budgets they converge; their earlier ``failure'' reflected the unequal protocol.
\actions Table~3.

\section{Reviewer 2}
\subsection*{Abstract, introduction, gap analysis, related work}
\begin{rcomment}{R2.1--R2.4}Highlight novelty and numbers in the abstract; increase motivation and recent literature; give explicit gap analysis; add a Related Work section citing HIC-net, the GLCM/LBP/LBGLCM/GLRLM/SFTA study, the CNN-Transformer EEG paper and the Jena-API paper.\end{rcomment}
\response The abstract now leads with the leakage finding and the corrected protocol, states the corrected headline numbers with CIs, and states plainly that the previously claimed components are within seed noise---the novelty is the audit and what it reveals, not the additions. The introduction is shortened, opens from the concrete problem, and ends with an explicit gap statement. A Related Work section is added with four subsections; HIC-net and the GLCM/LBP study are discussed substantively as the handcrafted-feature lineage our Gabor channel descends from; the EEG and clinical-text papers are cited briefly as deep learning in other clinical modalities. Recent references (Cellpose-SAM, SAM, MedSAM, StarDist, HoVer-Net) are added. \internalnote{Two bibliography entries (EEG, Jena API) carry \texttt{\textbackslash pending} markers for full bibliographic details---verify before submission.}

\subsection*{Figures}
\begin{rcomment}{R2.5}Figures should be drawn more effectively, supported with formulas, and show the novel pieces.\end{rcomment}
\response All figures are regenerated from code: leakage audit (Fig.~1), comparison with CIs and per-image distributions (Fig.~2), ablation (Fig.~3), real Gabor kernels/responses (Fig.~5), encoder ladder (Fig.~6), training curves (Fig.~7). Equations for the edge channel, projection, objective and sampler are in Methods.

\subsection*{Ablations, statistics, dataset, baselines}
\begin{rcomment}{R2.6--R2.9}Expand ablations; add significance tests; more dataset detail; more SOTA baselines.\end{rcomment}
\response Addressed under R1.6/R1.8/R1.17 (ablations: 11 single-change rows, 4 detectors, 6 sensitivity runs, 6-encoder ladder, vanilla anchor), R1.12 (statistics), R1.1/R1.2/R1.18 (dataset), and R1.7/R1.20 (baselines). Cellpose-SAM is added as the nucleus-specialist reference: zero-shot @@CELLPOSE_ZEROSHOT_DICE@@, fine-tuned @@CELLPOSE_FINETUNED_DICE@@ (precision @@CELLPOSE_FINETUNED_P@@). \internalnote{StarDist/HoVer-Net were not run: both depend on TensorFlow, which does not run on the RTX~5090 in this environment. Cellpose-SAM is the stronger and more current specialist reference; say so if asked.}

\section*{Closing}
The revised manuscript reports a lower headline number than the original, obtained under a protocol we can defend, and discloses the defects we found on the way. We believe it is a substantially more useful paper as a result.
\end{document}
